% ============================================================
%  Explicit gerbe class $[w_2^{\mathrm{per}}]$
%
%  Author: Raeez Lorgat.
% ============================================================

\providecommand{\Sp}{\mathrm{Sp}}
\providecommand{\Mp}{\mathrm{Mp}}
\providecommand{\Orth}{\mathrm{O}}
\providecommand{\SO}{\mathrm{SO}}
\providecommand{\Spin}{\mathrm{Spin}}
\providecommand{\Z}{\mathbb{Z}}
\providecommand{\Q}{\mathbb{Q}}
\providecommand{\R}{\mathbb{R}}
\providecommand{\C}{\mathbb{C}}
\providecommand{\A}{\mathbb{A}}
\providecommand{\CY}{\mathrm{CY}}
\providecommand{\NS}{\mathrm{NS}}
\providecommand{\Muk}{\widetilde{H}}
\providecommand{\disc}{\mathrm{disc}}
\providecommand{\Per}{\mathrm{Per}}
\providecommand{\Hilb}{\mathrm{Hilb}}
\providecommand{\Enr}{\mathrm{Enr}}
\providecommand{\Brau}{\mathrm{Br}}
\providecommand{\Num}{\mathrm{Num}}
\providecommand{\cA}{\mathcal{A}}
\providecommand{\cC}{\mathcal{C}}
\providecommand{\cL}{\mathcal{L}}
\providecommand{\cO}{\mathcal{O}}
\providecommand{\wper}{w_2^{\mathrm{per}}}
\providecommand{\kch}{\kappa_{\mathrm{ch}}}
\providecommand{\kcat}{\kappa_{\mathrm{cat}}}
\providecommand{\kBKM}{\kappa_{\mathrm{BKM}}}
\providecommand{\kfib}{\kappa_{\mathrm{fiber}}}

\section{The finite-discriminant period gerbe class}
\label{sec:wper-mukai-discriminant}

\subsection*{The obstruction class}

Let $L$ be an even integral lattice, $L^\vee$ its dual, and
\[
  A_L=L^\vee/L,\qquad q_L(x+L)=(x,x)\bmod 2\Z
\]
its discriminant quadratic module.  Let $\Gamma\subset\Orth(L)$ be
the monodromy group of a connected period component.  The
finite-discriminant part of the metaplectic period gerbe is the
pullback to $B\Gamma$ of the theta central extension
\[
  1\longrightarrow \{\pm1\}
  \longrightarrow \widetilde{\Gamma}_{q_L}
  \longrightarrow \Gamma
  \longrightarrow 1
\]
defined by the Weil representation of the finite quadratic module
$(A_L,q_L)$.  Its class is
\[
  e(L,\Gamma)\in H^2(\Gamma;\Z/2).
\]
For a period map $\Per_{\cC}:\pi_1(\cC)\to\Gamma$, the period-gerbe
class of $\cC$ is
\[
  [\wper(\cC)]_{\mathrm{disc}} =
  \Per_{\cC}^{*}e(L_\cC,\Gamma_\cC)
  \in H^2(\pi_1(\cC);\Z/2).
\]
If $A_L=0$, then the finite-quadratic-module contribution to
$[\wper]$ is zero.  If $A_L$ has a two-primary element $a$ with
$q_L(a)=\frac12\bmod\Z$, then the period gerbe is nontrivial on the
monodromy subgroup detecting $a$.

This formulation separates the orthogonal K3 and hyperkahler period
problem from the genus-two Siegel problem.  A pullback of the
Shimura-Weil cocycle on $\Sp_4$ requires a period
representation landing in $\Sp_4$, for example through a genus-two
Kuga-Satake or principally polarised abelian-surface window.

\subsection*{K3 surfaces}

Let $X$ be a projective K3 surface.  The full Mukai lattice of the
derived category is
\[
  \Muk(X,\Z)=H^0(X,\Z)\oplus H^2(X,\Z)\oplus H^4(X,\Z)
  \cong U^{\oplus4}\oplus E_8(-1)^{\oplus2},
\]
even unimodular of signature $(4,20)$ by Mukai's 1987 construction.
The ordinary weight-two period lattice
$H^2(X,\Z)\cong U^{\oplus3}\oplus E_8(-1)^{\oplus2}$ is also even
unimodular.  Therefore
\[
  A_{\Muk(X)}=0,\qquad A_{H^2(X)}=0,
\]
and the finite-quadratic-module class vanishes:
\[
  [\wper(D^b(X))]=0.
\]
The same conclusion holds at a CM point.  Complex multiplication
changes the Hodge endomorphism algebra of the period point while the
discriminant quadratic module remains zero.

\subsection*{Hilbert schemes of K3 type}

For a hyperkahler manifold $Y$ deformation-equivalent to
$\Hilb^n(K3)$, the Beauville-Bogomolov-Markman weight-two lattice is
\[
  H^2(Y,\Z)
  \cong U^{\oplus3}\oplus E_8(-1)^{\oplus2}
     \oplus \langle -2(n-1)\rangle .
\]
Thus
\[
  A_{H^2(Y)}\cong \Z/(2(n-1))\Z,\qquad
  q(\bar{1})=-\frac{1}{2(n-1)}\bmod 2\Z .
\]
Let $a=(n-1)\bar{1}$ be the unique element of order $2$ in
$A_{H^2(Y)}$.  Then
\[
  q(a)=-\frac{n-1}{2}\bmod 2\Z .
\]
Consequently the two-primary obstruction detected by the Markman
generator is
\[
  [\wper(Y)]_{\mathrm{disc}}=
  \begin{cases}
    0, & n\equiv 1 \pmod 2,\\
    \alpha_n, & n\equiv 0 \pmod 2,
  \end{cases}
\]
where $\alpha_n$ denotes the nonzero order-two character of the
monodromy action on the discriminant group.  This is a statement about
the hyperkahler weight-two period problem.  For $n>1$,
$D^b(\Hilb^n(K3))$ has CY dimension $2n$.

\subsection*{Enriques quotient surface}

Let $S$ be an Enriques surface and $\pi:X\to S$ its K3 double cover.
The free numerical lattice is
\[
  \Num(S)\cong U\oplus E_8(-1),
\]
and pullback to the invariant lattice of the K3 cover scales the
intersection form:
\[
  \pi^*\Num(S)\cong U(2)\oplus E_8(-2).
\]
Hence the invariant period lattice has two-elementary discriminant
group
\[
  A_{\pi^*\Num(S)}\cong(\Z/2\Z)^{10}.
\]
The canonical bundle $\omega_S$ supplies an additional order-two
torsion class.  The Enriques period gerbe is therefore nontrivial on
the quotient surface period stack.  It belongs to torsion-canonical
surface geometry.  The trivial-canonical CY$_2$ surface lane is the
separate K3 case above.

The associated Enriques metaplectic denominator lies in the
half-integral sector:
\[
  \Delta^{\Enr}_5\in S_{5/2}\bigl(\widetilde{K(2)},v_{\Enr}\bigr),
  \qquad
  \kBKM(\Delta^{\Enr}_5)=\frac52 .
\]
This row is separate from the primitive CHL order-$2$ denominator
$\Phi_2$, whose Borcherds weight is $\kBKM(\Phi_2)=4$.

\subsection*{Principally polarised abelian surfaces}

Let $(A,L)$ be a principally polarised abelian surface.  The full even
Mukai lattice
\[
  H^{\mathrm{even}}(A,\Z)=H^0(A,\Z)\oplus H^2(A,\Z)\oplus H^4(A,\Z)
  \cong U^{\oplus4}
\]
is unimodular.  Therefore the Mukai-discriminant contribution to the
period gerbe is zero:
\[
  A_{H^{\mathrm{even}}(A)}=0,\qquad
  [\wper(D^b(A))]_{\mathrm{Mukai}}=0.
\]

The theta-characteristic metaplectic cover over the Siegel modular
stack $\cA_2$ is a different class.  It is the central extension of
$\Sp_4(\Z)$ governing the Weil representation and theta multipliers.
It is detected by the Shimura-Weil and Rao cocycles on $\Sp_4(\A)$,
then pulled back along a genus-two period map.  It is separate from
the full Mukai-lattice discriminant class of $A$.

\subsection*{Shimura-Weil pullback}

Let $\rho:\pi_1(\cC)\to\Sp_4(\Z)$ be a period representation.  Howe's
adelic Weil representation and Rao's explicit cocycle define a central
extension class
\[
  [c_{\mathrm{Rao}}]\in H^2(\Sp_4(\Z);\Z/2).
\]
The Siegel metaplectic gerbe on $\cC$ is
\[
  [\wper^{\mathrm{Siegel}}(\cC)]
  =\rho^*[c_{\mathrm{Rao}}]\in H^2(\pi_1(\cC);\Z/2).
\]
For a K3-derived category with full Mukai monodromy in
$\Orth(\mathrm{II}_{4,20})$, this formula requires a specified
$\Sp_4$ realisation.  Absent that realisation, the relevant class is
$e(L,\Gamma)$ for the orthogonal finite quadratic module.

\subsection*{Computed cases}

\[
\begin{array}{l|l|l}
\text{period input} & \text{finite quadratic module} &
\text{period-gerbe conclusion} \\
\hline
D^b(K3) &
0 &
[\wper]=0 \\
\Hilb^n(K3)\text{ type} &
\Z/(2(n-1))\Z,\ q(1)=-1/(2(n-1)) &
[\wper]_{\mathrm{disc}}=(n-1)\bmod2 \\
\Enr &
(\Z/2\Z)^{10}\text{ on }U(2)\oplus E_8(-2),\ \omega_S\text{ torsion} &
[\wper]\neq0 \\
D^b(A),\ A\text{ abelian surface} &
0\text{ for }H^{\mathrm{even}}(A,\Z) &
[\wper]_{\mathrm{Mukai}}=0 \\
\cA_2\text{ Siegel stack} &
\text{theta multiplier on }\Sp_4(\Z) &
[\wper^{\mathrm{Siegel}}]=[c_{\mathrm{Rao}}]
\end{array}
\]

\subsection*{Borcherds and Mathieu normalisations}

The period-gerbe calculation uses the following Borcherds and
Mathieu normalisations as compatibility data.

For the primitive CHL family on $K3\times E$,
\[
  (c_1(0),c_2(0),c_3(0),c_4(0),c_6(0))=(10,8,6,4,2)
\]
and Borcherds' weight formula gives
\[
  (\kBKM(\Phi_1),\kBKM(\Phi_2),\kBKM(\Phi_3),
    \kBKM(\Phi_4),\kBKM(\Phi_6))=(5,4,3,2,1).
\]
At $N=1$, $\Phi_1=\Delta_5$ and $\mathrm{wt}(\Delta_5)=5$.
For the compact total space $K3\times E$ the four scalars
\[
  \kcat(K3\times E)=0,\qquad
  \kch^{\mathrm{Heis}}(K3\times E)=3,\qquad
  \kBKM(\Delta_5)=5,\qquad
  \kfib(K3)=24
\]
come from four distinct constructions.

The Gritsenko-Clery reflective catalogue is a separate eight-row
catalogue with twice-weight tuple
\[
  (10,4,6,2,4,1,3,2),
\]
equivalently weights
\[
  (5,2,3,1,2,\tfrac12,\tfrac32,1).
\]
It is indexed by reflective paramodular data.  The five CHL orders form
the separate primitive denominator row.

The ordinary Mathieu table has $26$ conjugacy classes.  The frame
shapes used here include the following corrected entries:
\[
\begin{array}{c|c}
\text{class} & \text{frame shape} \\
\hline
3B & 3^8 \\
4A & 2^4 4^4 \\
4B & 1^4 2^2 4^4 \\
6B & 6^4 \\
10A & 2^2 10^2 \\
12A & 2^1 4^1 6^1 12^1 \\
12B & 12^2
\end{array}
\]
These are ordinary $M_{24}$ frame shapes, separate from CHL
denominator weights and from the Gritsenko-Clery catalogue.

\paragraph{Primary literature.}
Giraud 1971, \emph{Cohomologie non abelienne}, III.2.1, for gerbe
classes; Shimura 1975, \emph{Annals of Mathematics} 102, for
metaplectic splitting in the Siegel setting; Howe 1979, \emph{Proc.
Symp. Pure Math.} 33, and Rao 1993, \emph{Pacific Journal of
Mathematics} 157, for the Weil cocycle; Nikulin 1979, \emph{Math.
USSR Izvestija} 14, for discriminant forms and Enriques lattices;
Mukai 1987, \emph{Inventiones Mathematicae} 77, for Mukai lattices;
Markman 2010, \emph{IMRN}, Theorem 1.2, for K3$^{[n]}$ monodromy and
the Beauville-Bogomolov lattice; Borcherds 1998, \emph{Inventiones
Mathematicae} 132, Theorem 13.3, for the weight formula
$\operatorname{wt}\Psi(f)=c(0)/2$; Gritsenko-Nikulin 1998 and
Gritsenko 1999 for $\Delta_5$ and the Enriques metaplectic
denominator; Gaberdiel-Hohenegger-Volpato 2010 and
Cheng-Duncan-Harvey 2014 for the Mathieu frame-shape normalisation.
